\begin{table}[htbp!]
\centering
\caption{F1-\textit{score} por categoria de consulta e por modelo --- referência em nuvem (Groq)}
\label{tab:resultados_por_categoria_groq}
\small
\ajustartabela{%
\begin{tabular}{|l|c|c|}
\hline
\textbf{Categoria} & \textbf{GPT-OSS 20B (Groq)} & \textbf{GPT-OSS 120B (Groq)} \\
\hline
Simples (n=100) & 0,985 & 0,979 \\
Compostas (n=180) & 0,937 & 0,950 \\
Código MI/INOM (n=57) & 0,985 & 0,986 \\
Tempo relativo (n=60) & 0,938 & 0,933 \\
Ordenação (n=29) & 0,837 & 0,926 \\
Ambíguas/informais (n=197) & 0,952 & 0,951 \\
\hline
\textbf{Média ponderada} & 0,946 & 0,956 \\
\hline
\end{tabular}}
\fonte{Elaborado pelos autores. Uma consulta pode pertencer a mais de uma categoria; n = consultas distintas da categoria nas métricas principais (cada uma pontuada em todas as repetições). As consultas fora do domínio (categoria F), cujo gabarito é não chamar a ferramenta, não têm campos e aparecem na Figura de acurácia por categoria.}
\end{table}
